\chapter{CONCLUSIONS}
\label{ch:conclusions}

The Modified State Observer has been extended to discrete time and its stability
established by Lyapunov analysis. \rev{The proof evaluates the first difference
of the Lyapunov function exactly, which exposes the cancellations that the
adaptation law was constructed to produce. The resulting conditions are
$\sigmax(F-\Km)<1$ together with an upper bound on the adaptation gain, both
satisfiable over an open set of designs, and both reducing to the expected
limiting behavior as the adaptation gain tends to zero. The analysis bounds the
state estimation error and the functional approximation error
$\Wt_k^T\phi_k$; boundedness of the weight estimate itself follows either from
persistency of excitation or, without that requirement, from
$\sigma$-modification of the update law.}

\rev{In simulation on the corrected robot model, the observer satisfies its
conditions with margin, settles inside the predicted ultimate bound, and under
deterministic parameter error identifies the uncertainty an order of magnitude
more accurately than an augmented-state Kalman filter. With realistic sensor
noise the advantage is confined to the velocity channel. The diagonal error
injection that makes the proof simplest cannot use the model to reject the
structured error of the accelerometer-derived tilt, which a Kalman-gain
injection recovers empirically.}

\rev{A command-filtered neural backstepping controller has been developed for
the two-wheeled inverted pendulum, an underactuated system with unmatched
uncertainty. The design assigns damping to every error coordinate, retains and
explicitly cancels the coupling between the position and attitude error
subsystems, and uses command filtering so that each network approximates a
function of measured state and reference alone. Its boundedness result is
conditional on the command-filter error assumption. For the TWIP that
assumption does not hold: tracking the non-minimum-phase position output
places a closed-loop pole at the right-half-plane zero of $x_1(s)/u(s)$, and the
design could not be stabilized. The revision~7 two-step extra control does
balance the robot in every case, with mixed effect on tracking accuracy.}

Four directions follow naturally from this work.

\paragraph{A realizable backstepping output.} \rev{Repeating the four-step
construction on the flat output $y=Cx$, $CB=CAB=CA^2B=0$, which has no zero
dynamics, or restructuring it as a cascade with a position loop slower than the
zero, would retain the damping and command-filtering ideas of
Chapter~\ref{ch:control} while removing the obstruction identified in
Section~\ref{sec:cfinstab}.}

\paragraph{Output feedback.} The present development assumes full state
measurement, so the DMSO is properly a filter and uncertainty identifier. The
structure of Proposition~\ref{prop:exact} carries over to $y_k=Cx_k$ with
$A=F-\Km C$, and the extension requires only a detectability assumption. This
is the most direct route to a result that applies to systems where the tilt rate
or the wheel velocity is not directly measured.

\paragraph{Nonsmooth input nonlinearities.} Backlash violates the smoothness
assumption underlying the universal approximation property, and the residual
limit cycles observed both in simulation and on hardware are its signature. A
network cannot cancel what it cannot approximate. Two remedies are worth
pursuing: an explicit deadzone- and backlash-inverse placed ahead of the plant,
or a discontinuous term in the control law that is robust to the deadband rather
than attempting to estimate it.

\paragraph{Deterministic timing.} \rev{The achievable sample rate on the
platform was set by telemetry rather than by the control computation, and the
encoder velocity quantization scales inversely with the sample period. Both are
implementation choices rather than properties of the algorithm, and both bound
what any estimator can deliver. Appendix~\ref{app:validation} treats the loop
rate as a design variable and gives the trade.}
